\section{A support-free step and the least contrary index}
\label{ms:smoothing}

The role of smoothing here is modest but useful. A sharp preceding index
forces any next contrary score into the narrow resource regime used by the
local analysis. The estimate applies to every standardized marginal with a
finite third absolute moment; it does not use contact or support acquisition.

\begin{lemma}[Elementary Gaussian constants]
\label{ms:gaussian-constants}
For the standard normal density $\phi$,
\[
 \|\phi''\|_\infty=\phi(0)<\frac25,
 \qquad \|\phi'''\|_\infty\le\frac2\pi,
 \qquad \frac1{4\pi}<\frac2{25}.
\]
Moreover $819/2000<\CE<20487/50000<41/100$ and
$\phi(0)>0.39886>0.3988=997/2500$.
\end{lemma}
\begin{proof}
Since $\phi''(x)=(x^2-1)\phi(x)$, the norm on $x^2\le1$ is at most
$\phi(0)$. On $x^2\ge1$, the maximum of
$(x^2-1)e^{-x^2/2}$ is $2e^{-3/2}<1$, attained when $x^2=3$.
Thus the norm is exactly $\phi(0)$.

For completeness, the Gaussian transform can be obtained without an
inversion theorem for the distribution being studied. Differentiating
$J(t)=\int_\R e^{-x^2/2}e^{itx}\,dx$ under the integral and integrating by
parts give $J'(t)=-tJ(t)$. The Gaussian integral, evaluated by squaring it
and using polar coordinates, gives $J(0)=\sqrt{2\pi}$. Hence
$J(t)=\sqrt{2\pi}e^{-t^2/2}$; applying the same formula with the variables
interchanged gives
\[
 \phi(x)=\frac1{2\pi}\int_\R e^{-t^2/2}e^{-itx}\,dt.
\]
Three differentiations are justified by the integrability of
$|t|^3e^{-t^2/2}$. Consequently
\[
 \|\phi'''\|_\infty
 \le\frac1{2\pi}\int_\R|t|^3e^{-t^2/2}\,dt=\frac2\pi.
\]
The last integral equals four after the substitution $w=t^2/2$ on the
positive half-line.

Here are elementary rational comparisons for all numerical constants used
in this section. Put $a=\arctan(1/5)$ and $b=\arctan(1/239)$.
Tangent addition gives $\tan(2a)=5/12$, $\tan(4a)=120/119$ and
$\tan(4a-b)=1$. Also
\[
 0<4(1/5-1/(3\cdot5^3))-1/239<4a-b<4/5<1<\pi/2.
\]
The arctangent inequalities follow by integrating
$1-t^2<1/(1+t^2)<1$, while $\pi/2=2\arctan(1)>1$ follows by integrating
$1/(1+t^2)>1/2$ on $[0,1)$. It follows that $\pi=16a-4b$ and
\[
 \pi>16(1/5-1/(3\cdot5^3))-4/239
      =281476/89625>25/8.
\]
This proves $\phi(0)<2/5$ and $1/(4\pi)<2/25$.

To bound $\CE$ more closely, define
\[
 s_k(x)=\sum_{j=0}^{k-1}\frac{(-1)^jx^{2j+1}}{2j+1},\quad
 \pi_-=16s_6(1/5)-4s_1(1/239),\quad
 \pi_+=16s_5(1/5)-4s_2(1/239).
\]
The remainder in the finite geometric expansion of $1/(1+t^2)$ has its
stated alternating sign, so integration gives $\pi_-<\pi<\pi_+$.
Let $r_-=3162277660/10^9$ and $r_+=3162277661/10^9$. Squaring shows
$r_-^2<10<r_+^2$. The following comparisons are exact rational ones:
\[
 \left(\frac{819}{2000}\right)^2
 <\frac{19+6r_-}{72\pi_+}
 <\CE^2=\frac{19+6\sqrt{10}}{72\pi}
 <\frac{19+6r_+}{72\pi_-}
 <\left(\frac{20487}{50000}\right)^2.
\]
For example, each outer comparison follows by multiplying by its positive
rational denominator; the displayed finite sums specify those denominators
completely. All quantities are positive, so taking square roots proves the
claimed bounds. In particular, $0.4095<\CE<0.40974<0.41$.
The same upper endpoint for $\pi$ is smaller than $22/7$.
Since $(0.39886)^2<7/44$, we also obtain
$\phi(0)>0.39886>0.3988$.
\end{proof}

\begin{theorem}[One-step smoothing]
\label{ms:one-step}
For every integer $n\ge2$,
\begin{equation}
 C_n\le\sqrt{\frac n{n-1}}C_{n-1}
       +\frac{\sqrt n}{15(n-1)^{3/2}}
       +\frac{2\sqrt n}{25(n-1)^2}.
\label{ms:smoothing-recurrence}
\end{equation}
\end{theorem}
\begin{proof}
Fix $P\in\mathcal P$ and a raw threshold $H$. Write $v=n-1$ and
$G_v(x)=\Phi(x/\sqrt v)$. Conditioning on the last summand gives
\[
 \left|F_{n,P}(H)-\E G_v(H-X_n)\right|
 \le C_v\rho/\sqrt v.
\]
This uses the defining bound for the same marginal at the preceding index;
the CDF of the sum is not differentiated.

Taylor's theorem, applied only to the smooth function $G_v$, gives
\[
 \left|G_v(H-y)-G_v(H)+yG_v'(H)-\tfrac12y^2G_v''(H)\right|
 \le\frac{\phi(0)}{6v^{3/2}}|y|^3.
\]
Taking expectations and using the first two moments yields
\[
 \left|\E G_v(H-X)-G_v(H)-\tfrac12G_v''(H)\right|
 \le\frac{\phi(0)\rho}{6v^{3/2}}<\frac\rho{15v^{3/2}}.
\]
For the change in Gaussian variance, direct differentiation gives
$\partial_vG_v=G_v''/2$ and $\partial_v^2G_v=G_v^{(4)}/4$.
Taylor's integral remainder in the variance therefore gives
\[
 G_{v+1}(H)-G_v(H)-\tfrac12G_v''(H)
   =\frac14\int_0^1(1-t)G_{v+t}^{(4)}(H)\,dt.
\]
As $G_w^{(4)}(H)=w^{-2}\phi'''(H/\sqrt w)$, the variance remainder
has absolute value at most
\[
 \frac{\|\phi'''\|_\infty}{8v^2}
       \le\frac1{4\pi v^2}<\frac2{25v^2}.
\]
Combine the three bounds, multiply by $\sqrt n/\rho$, and use $\rho\ge1$.
The resulting estimate holds for every $P,H$, so taking both suprema proves
\eqref{ms:smoothing-recurrence}.

Only a constant multiple of $|X|^3$ was integrated in the non-Gaussian
remainder. Thus rare remote atoms and unbounded support cause no additional
assumption. The bound is uniform in $H$ and controls both signs. Strict
events follow as limits of increasing closed thresholds, with the same
normal value at the limit.
\end{proof}

\begin{lemma}[Entry at a sharp predecessor]
\label{ms:first-failure-entry}
If $n\ge42$ and $C_{n-1}\le\CE\le C_n$, then
\[
 0\le n(C_n-\CE)<\frac{3003}{10250}<\frac3{10},\qquad
 C_n<\frac{73}{175}<\frac{2113}{5000},\qquad
 n(C_{n-1}-C_n)_+=0.
\]
For every marginal $P\in\mathcal P$,
$n(R_{n-1}(P)-C_n)_+=0$ as well.
\end{lemma}
\begin{proof}
Let $s=\sqrt{n/(n-1)}$ and $c=\CE$. From
\eqref{ms:smoothing-recurrence} and
$n(s-1)=(n/(n-1))/(s+1)$,
\[
 n(C_n-c)\le\frac n{n-1}
 \left\{\frac c{s+1}+\frac s{15}
                         +\frac{2s}{25\sqrt{n-1}}\right\}.
\]
For $n\ge42$, the ratio $n/(n-1)$ is at most $42/41$,
$\sqrt{n-1}>6$, and $1<s<81/80$. The latter follows at the endpoint
from $41\cdot81^2-42\cdot80^2=201>0$ and then from monotonicity.
Using $c<41/100$ gives the strict bound
\[
 n(C_n-c)<\frac{42}{41}
 \left(\frac{41}{200}+\frac{27}{400}+\frac{27}{2000}\right)
 =\frac{3003}{10250}.
\]
The reserve below $3/10$ is $36/5125$. Consequently
$C_n<c+0.3/n<0.41+0.3/42=73/175$, whose reserve below $0.4226$ is
$191/35000$. Finally $R_{n-1}(P)\le C_{n-1}\le c\le C_n$, so both
positive parts vanish.
\end{proof}

\begin{proposition}[Equality-inclusive least contrary index]
\label{ms:least-contrary}
Assume $C_j\le\CE$ for $1\le j\le41$. If some $C_n\ge\CE$ with
$n\ge42$, there is a least such index, and at that index
Lemma~\ref{ms:first-failure-entry} holds, while
\[
 C_{n-1},C_{n-2}\le\CE<0.41<0.4253.
\]
Thus any independently established implication excluding an attained
positive closed witness under these entry conditions excludes that least
index. If such an implication is valid for every actual $n\ge42$, it gives
$C_n<\CE$ throughout that range.
\end{proposition}
\begin{proof}
The positive integers are well ordered. Below the least chosen index,
the indices from $42$ onward have $C_j<\CE$, and the earlier ones have the
assumed weak bound. At $n=42$ the required predecessors are $41$ and $40$;
at later indices the same conclusion holds. Apply the scalar lemma and
then Lemma~\ref{ms:attainment}. Every premise of the stated independent
exclusion implication is now available, including equality at the current
score. The contradiction proves the conclusion.
\end{proof}

Choosing the least index with $C_n>\CE$ would give only a weak conclusion.
Neither the sharp finite base nor the all-index exclusion implication is
proved by the least-index argument itself.
